% §8.2 Channel Capacity in the BSC
%
% Cross-reference resolution status:
%
% DEFINED here:
%   \label{sec:bsc-capacity}                      -- this section
%   \label{prop:bsc-capacity}                     -- the BSC capacity result
%
% ALREADY DEFINED (chapter 2):
%   \Cref{sec:bsc}                                -- chapter 2 §2.1
%   \Cref{def:bsc}                                -- chapter 2 §2.1 BSC definition
%
% ALREADY DEFINED (chapter 8 §8.1, Plan 8 Task 2):
%   \Cref{sec:entropy-refresher}                  -- §8.1
%   \eqref{eq:binary-entropy}                     -- §8.1 binary entropy function
%   \eqref{eq:mutual-info}                        -- §8.1 mutual information
%   \eqref{eq:capacity-def}                       -- §8.1 capacity definition
%
% FORWARD REFS within chapter 8 (resolve in later Plan 8 tasks):
%   \Cref{sec:entropy-floor}                      -- §8.3 (Task 4)
%
% §8.2 applies the information-theoretic tools of §8.1 to the BSC of
% chapter 2 and derives the closed-form capacity $C = 1 - H(\fprate)$,
% achieved at the uniform input distribution. The section recovers the
% endpoints (one bit per use at $\fprate = 0$, zero bits per use at
% $\fprate = 1/2$) and bridges forward to the source-coding-theorem
% reading that §8.3 applies to approximate-set representations.

\section{Channel Capacity in the BSC}
\label{sec:bsc-capacity}

\subsection*{Recall: the BSC}

The binary symmetric channel of \Cref{sec:bsc} takes a single Boolean
input $X \in \{\top, \bot\}$ and produces a Boolean output $Y \in
\{\top, \bot\}$. \Cref{def:bsc} fixes a single parameter $\fprate \in
(0, 1)$ that governs the channel: each transmitted bit is reproduced
faithfully with probability $1 - \fprate$ and flipped to the opposite
value with probability $\fprate$, the symmetry being that the flip
probability does not depend on which value was sent. The channel
matrix
\[
  Q \;=\;
  \begin{pmatrix}
    1 - \fprate & \fprate \\
    \fprate & 1 - \fprate
  \end{pmatrix}
\]
records this: rows indexed by the input, columns by the output, both
diagonal entries equal to $1 - \fprate$ (correct transmission) and
both off-diagonal entries equal to $\fprate$ (flip). The BSC is the
common abstraction behind chapter~1's biased coin, noisy
communication line, and binary classifier; capacity is its
quantitative information-theoretic summary.

\subsection*{Capacity, specialized to the BSC}

The channel capacity defined in \eqref{eq:capacity-def} of
\Cref{sec:entropy-refresher} is
\[
  C \;=\; \max_{P_X} I(X; Y),
\]
the largest mutual information achievable by choosing the input
distribution $P_X$ on $\{\top, \bot\}$ freely while the conditional
distribution $\Pr[Y \mid X]$ is held fixed by the channel matrix $Q$.
The maximization has a single scalar free parameter: write
$P_X(\top) = q \in [0, 1]$ and $P_X(\bot) = 1 - q$. Symmetry of the
BSC will collapse the optimum to $q = 1/2$, the uniform input. The
argument is short enough to write out in full.

\subsection*{Derivation}

We compute $I(X; Y) = H(Y) - H(Y \mid X)$, the second of the two
equivalent forms in \eqref{eq:mutual-info}, because $H(Y \mid X)$
turns out to be independent of the input distribution and the
optimization then reduces to maximizing $H(Y)$.

\medskip
\noindent\textbf{Step 1: $H(Y \mid X) = H(\fprate)$.}
Given input $X = \top$, the output $Y$ is a Bernoulli random
variable: $Y = \top$ with probability $1 - \fprate$ (correct
transmission) and $Y = \bot$ with probability $\fprate$ (flip). Its
entropy, by \eqref{eq:binary-entropy}, is $H(\fprate)$. The same
calculation with $X = \bot$ yields the same $H(\fprate)$, because the
two rows of $Q$ are permutations of each other. Averaging over the
input distribution then gives
\[
  H(Y \mid X)
  \;=\; q \cdot H(\fprate) + (1 - q) \cdot H(\fprate)
  \;=\; H(\fprate),
\]
a constant independent of $q$. The conditional entropy is the
intrinsic flip uncertainty of the channel and has nothing to do with
what the sender chooses to put in.

\medskip
\noindent\textbf{Step 2: $H(Y) \le 1$, with equality at $q = 1/2$.}
The output marginal probability $\Pr[Y = \top]$ is
\[
  \Pr[Y = \top]
  \;=\; q (1 - \fprate) + (1 - q) \fprate
  \;=\; \fprate + q (1 - 2 \fprate),
\]
which at $q = 1/2$ equals $\fprate + \tfrac{1}{2}(1 - 2 \fprate) =
\tfrac{1}{2}$. The output is uniform on $\{\top, \bot\}$, and its
entropy is the maximum value of the binary entropy function, $H(1/2)
= 1$ bit. For any other $q$, $\Pr[Y = \top]$ is some value in $[0,
1]$ different from $\tfrac{1}{2}$, and $H(Y) < 1$ by strict concavity
of $H$ with maximum at $1/2$.

\medskip
\noindent\textbf{Step 3: combine.}
The mutual information is
\[
  I(X; Y)
  \;=\; H(Y) - H(Y \mid X)
  \;=\; H\bigl(\fprate + q(1 - 2 \fprate)\bigr) - H(\fprate),
\]
with the second term constant. Maximizing over $q$ pushes the first
term to its ceiling of $1$, attained at $q = 1/2$. The capacity is
the maximum:
\[
  C \;=\; 1 - H(\fprate).
\]

\begin{proposition}[BSC capacity]
\label{prop:bsc-capacity}
The capacity of the binary symmetric channel of \Cref{def:bsc} with
crossover probability $\fprate \in (0, 1)$ is
\begin{equation}
\label{eq:bsc-capacity}
  C \;=\; 1 - H(\fprate)
  \;=\; 1 + \fprate \log_2 \fprate
       + (1 - \fprate) \log_2 (1 - \fprate)
  \quad \text{bits per channel use,}
\end{equation}
attained by the uniform input distribution $P_X(\top) = P_X(\bot) =
1/2$.
\end{proposition}

A single number summarizes the BSC's quantitative information
content. Once $\fprate$ is fixed, the channel admits at most $C(\fprate)$
bits per use, no matter what coding scheme the sender devises.
Shannon's noisy-channel coding theorem then promises that this
ceiling is asymptotically achievable: for any rate $R < C$ and any
$\varepsilon > 0$, there exist block codes of large enough length
whose decoding error rate is below $\varepsilon$. The BSC capacity
$1 - H(\fprate)$ is thus both an upper bound on what is possible and
the asymptotic rate that good codes actually achieve.

\subsection*{Reading the result}

The closed form $C = 1 - H(\fprate)$ has a sharp shape worth fixing
in mind, because §8.3 uses the same binary entropy function in a
structurally similar role for approximate sets.

\medskip
\noindent\textbf{$\fprate = 0$: capacity $1$.}
A noiseless channel transmits one bit per use without error. Setting
$\fprate = 0$ in \eqref{eq:bsc-capacity} gives $H(0) = 0$, so
$C = 1 - 0 = 1$. The output equals the input deterministically; one
bit goes in, one bit comes out.

\medskip
\noindent\textbf{$\fprate = 1/2$: capacity $0$.}
A maximally noisy channel transmits zero bits per use. Setting
$\fprate = 1/2$ gives $H(1/2) = 1$, so $C = 1 - 1 = 0$. The output
is a fair coin flip, statistically independent of the input, and no
amount of clever coding extracts information from it.

\medskip
\noindent\textbf{$\fprate = 1$: capacity $1$ again.}
A channel that always flips the bit is just as informative as a
noiseless one: the receiver applies the deterministic relabeling
$\top \leftrightarrow \bot$ and recovers the input perfectly. The
formula reflects this: $H(1) = 0$, so $C = 1$. The function $C(\fprate) =
1 - H(\fprate)$ is symmetric about $\fprate = 1/2$, mirroring the
symmetry of the binary entropy function.

\medskip
\noindent\textbf{Intermediate values.}
At $\fprate = 0.01$, $H(0.01) \approx 0.081$, so $C \approx 0.919$
bits per use: a $1\%$-error channel delivers roughly $92\%$ of one
bit per transmission, a small but real penalty. At $\fprate = 0.1$,
$C \approx 1 - 0.469 = 0.531$ bits per use, slightly more than half
a bit. The shape of $C$ on $[0, 1/2]$ is steep near $\fprate = 0$
(the first bit of crossover noise costs little) and shallow near
$\fprate = 1/2$ (the last bits of noise cost almost everything).
The asymmetry has a direct operational reading. A channel with $1\%$
crossover is almost as good as a noiseless one for practical coding
purposes; a channel with $40\%$ crossover, by contrast, has $C \approx
0.029$ bits per use, and getting any useful throughput out of it
requires a very long block code spreading each source bit over many
channel uses. The chapter~8 notebook in
\texttt{python/notebooks/ch08-information-bounds.ipynb} plots the
full curve and tabulates representative values.

\subsection*{From channel capacity to representation cost}

\Cref{prop:bsc-capacity} is the canonical first application of
\Cref{sec:entropy-refresher}'s tools and the quantitative completion
of chapter~2's BSC. The same machinery underlies a different but
related question that occupies the rest of this chapter: how many
bits are needed to \emph{represent} an approximate set with target
false-positive rate $\fprate$? Shannon's source-coding theorem,
mentioned briefly in \Cref{sec:entropy-refresher}, gives the relevant
form. Any prefix-free binary encoding of a random source emitting
$n$ symbols of entropy $H$ each requires on average at least $n H$
bits. \Cref{sec:entropy-floor} applies this lower bound to the bit
string that an approximate-set data structure stores: the bits in
that string have a constrained distribution governed by $\fprate$,
and the resulting per-element entropy gives a floor of $-\log_2
\fprate$ bits per stored element for any FNR-zero representation.
The BSC capacity derived here and the entropy floor derived there
are two facets of the same source-coding principle, with the role
of the binary entropy function $H$ holding throughout.
